% Standalone problem document, generated from the adjacent README.md.
% Compile with XeLaTeX. Mathematical content is maintained in Markdown.
\documentclass[11pt,a4paper]{article}
\usepackage[fontsize=10.5pt]{scrextend}
\usepackage[margin=22mm,headheight=15pt,headsep=9mm,footskip=11mm]{geometry}
\usepackage{fontspec}
\setmainfont{texgyrepagella}[Extension=.otf,UprightFont=*-regular,BoldFont=*-bold,ItalicFont=*-italic,BoldItalicFont=*-bolditalic]
\setsansfont{texgyreheros}[Extension=.otf,UprightFont=*-regular,BoldFont=*-bold,ItalicFont=*-italic,BoldItalicFont=*-bolditalic]
\setmonofont{texgyrecursor}[Extension=.otf,UprightFont=*-regular,BoldFont=*-bold,ItalicFont=*-italic,BoldItalicFont=*-bolditalic,Scale=MatchLowercase]
\usepackage{amsmath,amssymb}
\usepackage{unicode-math}
\setmathfont{texgyrepagella-math.otf}
\usepackage{xcolor,microtype,enumitem,needspace,fancyhdr}
\usepackage{bookmark}
\definecolor{ink}{HTML}{17344B}
\definecolor{accent}{HTML}{197D80}
\definecolor{muted}{HTML}{596773}
\hypersetup{colorlinks=true,linkcolor=accent,urlcolor=accent,pdftitle={FR-04 - Universal
exponential deterioration of minimally redundant real phase-retrieval
frames},pdfauthor={Open Problems in Numerical Linear Algebra}}
\urlstyle{same}
\setlength{\parindent}{0pt}
\setlength{\parskip}{5pt plus 1pt minus 1pt}
\linespread{1.04}
\setlength{\emergencystretch}{3em}
\widowpenalty=10000
\clubpenalty=10000
\displaywidowpenalty=10000
\allowdisplaybreaks[1]
\setlist{nosep,leftmargin=1.5em}
\providecommand{\tightlist}{\setlength{\itemsep}{0pt}\setlength{\parskip}{0pt}}
\providecommand{\pandocbounded}[1]{#1}
\pagestyle{fancy}
\fancyhf{}
\fancyhead[L]{\sffamily\footnotesize\color{muted}OPEN PROBLEMS IN NUMERICAL LINEAR ALGEBRA}
\fancyhead[R]{\sffamily\bfseries\footnotesize\color{accent}FR-04}
\fancyfoot[L]{\parbox[b]{0.9\textwidth}{\sffamily\fontsize{8}{10}\selectfont\color{muted}Editorial ratings; a bounded search cannot certify that a problem remains unsolved.\\Literature check: 2026-09-10}}
\fancyfoot[R]{\sffamily\footnotesize\color{muted}\thepage}
\renewcommand{\headrulewidth}{0.3pt}
\renewcommand{\headrule}{\hbox to\headwidth{\color{accent}\leaders\hrule height \headrulewidth\hfill}}
\makeatletter
\renewcommand{\section}{\Needspace{5\baselineskip}\@startsection{section}{1}{0pt}{12pt plus 2pt minus 2pt}{4pt}{\sffamily\bfseries\large\color{ink}}}
\renewcommand{\subsection}{\Needspace{4\baselineskip}\@startsection{subsection}{2}{0pt}{9pt plus 1pt minus 1pt}{3pt}{\sffamily\bfseries\normalsize\color{ink}}}
\makeatother
\setcounter{secnumdepth}{0}
\begin{document}
{\sffamily\small\color{muted}Frames and matrix designs\par}
\vspace{3pt}
{\raggedright\hyphenpenalty=10000\exhyphenpenalty=10000\sffamily\bfseries\fontsize{21}{24}\selectfont\color{ink}Universal
exponential deterioration of minimally redundant real phase-retrieval
frames\par}
\vspace{7pt}
{\sffamily\small\textbf{Difficulty:} challenging\quad\textbf{Importance:} interesting
to the community\par}
{\sffamily\small\bfseries\color{accent}Status: Open\par}
\vspace{4pt}
{\color{accent}\hrule height 0.8pt}
\vspace{6pt}
\textbf{Rating rationale:} A universal bound over every full-spark
measurement design requires a structural obstruction beyond the Gaussian
analysis; it matters for stable phase retrieval.

For \(n\geq2\), let \(A\in\mathbb R^{(2n-1)\times n}\) be full spark:
every choice of \(n\) rows is linearly independent. Write \(a_i^T\) for
its rows and define \[
L(A)=\max_{1\leq i\leq2n-1}\|a_i\|_2,\qquad
\omega(A)=\min_{\substack{T\subseteq\{1,\ldots,2n-1\}\\|T|=n}}
\sigma_{\min}(A_T).
\] Do absolute constants \(C>0\) and \(0<\beta<1\) exist such that \[
\omega(A)\leq C L(A)\beta^n
\] for every \(n\geq2\) and every such \(A\)?

For a full-spark matrix at this row count, the displayed definition is
equivalent to the Balan--Wang definition taking the minimum over row
sets whose complements fail to span \(\mathbb R^n\). It measures a
worst-conditioned square subproblem in inversion of the phaseless map
\(x\mapsto |Ax|\). The conjecture says that clever deterministic
measurement design cannot avoid exponential deterioration at minimal
real redundancy.

\subsection{References}\label{references}

\begin{enumerate}
\def\labelenumi{\arabic{enumi}.}
\tightlist
\item
  R. Balan and Y. Wang, \emph{Invertibility and robustness of phaseless
  reconstruction}, Applied and Computational Harmonic Analysis 38
  (2015), pp.~469--488. Conjecture 5.1 on printed p.~484; definitions
  and stability comparisons in Sections 3--5.
  \href{https://www.cscamm.umd.edu/publications/ACHA2015paper_CS-16-05.pdf}{Author-hosted
  paper}.
\item
  A. S. Bandeira et al., \emph{Randomstrasse 101: Open Problems of
  2025}, arXiv:2603.29571 (2026), Conjecture 20.
  \href{https://arxiv.org/html/2603.29571v1}{Paper}.
\item
  Y. Shmalo, \emph{Extreme least singular values of Gaussian row
  submatrices and a phase retrieval stability problem}, arXiv:2607.06249
  (2026). Abstract and Section 1.2/Corollary 1.2 give the critical real
  Gaussian asymptotic. \href{https://arxiv.org/abs/2607.06249}{Paper}.
\end{enumerate}

\subsection{Status check --- 2026-09-10}\label{status-check-2026-09-10}

Searched ``Balan Wang exponential conjecture solved'', ``phase retrieval
stability omega 2026'', and checked the latest abstract of
arXiv:2607.06249. The July 2026 result proves
\(\omega(A_n)=4^{-n+o_P(n)}\) for independent standard Gaussian entries.
A probabilistic result for this ensemble does not prove the uniform
inequality over all full-spark matrices. The separate Gaussian-base
question should therefore be excluded, while this deterministic
conjecture remains open in the screened literature.

\textbf{Audit update (2026-09-10):} Rechecked Conjecture 20 and Shmalo's
July 2026 Gaussian asymptotic, then searched for universal Balan--Wang
bounds. A high-probability ensemble theorem is not a deterministic bound
for every full-spark matrix, so the target remains open. This is a
bounded literature check, not a proof that no solution exists.
\end{document}
